\documentclass[12pt]{article}
\usepackage{amsmath,amssymb,amsthm}
\usepackage[margin=1in]{geometry}

\newtheorem{theorem}{Theorem}
\newtheorem{lemma}[theorem]{Lemma}

\title{Problem 191: Prize Strings}
\author{Project Euler}
\date{}

\begin{document}
\maketitle

\section{Problem Statement}

Over a 30-day period, how many strings of length 30 over the alphabet $\{O, A, L\}$ satisfy:
\begin{enumerate}
    \item At most one $L$ appears in the string.
    \item No three consecutive $A$'s appear.
\end{enumerate}

\section{Mathematical Foundation}

Let $\Sigma = \{O, A, L\}$ and let $\mathcal{P}_n \subseteq \Sigma^n$ denote the set of prize strings of length~$n$.

\begin{theorem}
Let $g(n)$ denote the number of strings of length $n$ over $\{O, A\}$ containing no three consecutive $A$'s. Then
\[
    g(n) = g(n-1) + g(n-2) + g(n-3), \quad n \ge 3,
\]
with initial values $g(0) = 1$, $g(1) = 2$, $g(2) = 4$.
\end{theorem}

\begin{proof}
Partition the valid strings of length $n$ by their suffix:
\begin{itemize}
    \item Strings ending in $O$: the prefix of length $n-1$ is any valid string, contributing $g(n-1)$.
    \item Strings ending in $AO$: the prefix of length $n-2$ is any valid string, contributing $g(n-2)$.
    \item Strings ending in $AA$ whose character at position $n-3$ is $O$ (since $AAA$ is forbidden): the prefix of length $n-3$ is any valid string, contributing $g(n-3)$.
\end{itemize}
These cases are exhaustive and mutually exclusive. The base cases are verified by enumeration: $g(0) = 1$ (empty string), $g(1) = 2$ ($\{O, A\}$), $g(2) = 4$ ($\{OO, OA, AO, AA\}$). \qed
\end{proof}

\begin{theorem}
The total number of prize strings of length $n$ is
\[
    |\mathcal{P}_n| = g(n) + \sum_{k=0}^{n-1} g(k) \cdot g(n-1-k).
\]
\end{theorem}

\begin{proof}
A prize string either contains zero $L$'s or exactly one $L$.

\textbf{Case 1} (zero $L$'s): The string is over $\{O, A\}$ with no $AAA$, counted by $g(n)$.

\textbf{Case 2} (exactly one $L$ at position $k$): The prefix of length $k$ is a valid no-$L$ string ($g(k)$ choices), and the suffix of length $n-1-k$ is an independent valid no-$L$ string ($g(n-1-k)$ choices). The $L$ resets the consecutive-$A$ counter, making the prefix and suffix constraints independent. Summing over $k \in \{0, \ldots, n-1\}$ gives $\sum_{k=0}^{n-1} g(k) \cdot g(n-1-k)$. \qed
\end{proof}

\begin{lemma}
$g(30) = 53798080$, computable in $O(n)$ arithmetic operations.
\end{lemma}

\begin{proof}
Direct application of the tribonacci-like recurrence with a sliding window. \qed
\end{proof}

\section{Algorithm}

\begin{verbatim}
function count_prize_strings(n):
    g = array of size n+1
    g[0] = 1; g[1] = 2; g[2] = 4
    for i from 3 to n:
        g[i] = g[i-1] + g[i-2] + g[i-3]
    total = g[n]
    for k from 0 to n-1:
        total += g[k] * g[n-1-k]
    return total
\end{verbatim}

\section{Complexity Analysis}

\begin{itemize}
    \item \textbf{Time:} $O(n)$ for the recurrence and convolution sum.
    \item \textbf{Space:} $O(n)$ to store $g(0), \ldots, g(n)$. Reducible to $O(1)$ with a two-pointer accumulation.
\end{itemize}

\section{Answer}

\[
\boxed{1918080160}
\]

\end{document}
